% Currículo em português. Gera assets/cv/thales-amaral-cv-pt.pdf (ver build.sh)
\documentclass[10pt,a4paper]{article}
% Preâmbulo comum aos currículos PT e EN. Compilar com ./build.sh

\usepackage[utf8]{inputenc}
\usepackage[T1]{fontenc}
\usepackage{lmodern}
\usepackage[margin=1.2cm,top=0.8cm,bottom=0.8cm]{geometry}
\usepackage[hidelinks]{hyperref}
\usepackage{enumitem}
\usepackage{titlesec}

% ATS: texto do PDF extraível com acentos e ligaduras corretos
\input{glyphtounicode}
\pdfgentounicode=1

\pagestyle{empty}

\setlength{\parindent}{0pt}
\setlength{\parskip}{0pt}

\setlist[itemize]{
    leftmargin=1.2em,
    itemsep=0.5pt,
    topsep=0.5pt,
    parsep=0pt
}

\titleformat{\section}
    {\normalsize\bfseries}
    {}{0pt}{}
    [\vspace{-4pt}\rule{\textwidth}{0.5pt}]

\titlespacing*{\section}{0pt}{4pt}{2pt}

% \entry{cargo/título}{datas}{instituição}{local}
\newcommand{\entry}[4]{
    \vspace{-2pt}
    \textbf{#1} \hfill \textbf{#2}\\
    \textit{#3} \hfill #4\\[-4pt]
}

% ATS: não hifenizar (palavras partidas como "Whis-per" atrapalham a busca por palavras-chave)
\hyphenpenalty=10000
\exhyphenpenalty=10000
\emergencystretch=2em

\hypersetup{
    pdftitle={Thales Amaral - Currículo},
    pdfauthor={Thales Amaral},
    pdfsubject={Engenheiro de IA e Dados},
    pdfkeywords={IA, LLM, RAG, MCP, Python, FastAPI, Airflow, Engenharia de Dados, Machine Learning}
}

\begin{document}

% ================================================================
% CABEÇALHO
% ================================================================

\begin{center}
    {\Large\bfseries Thales A. Amaral}\\[2pt]
    Engenheiro de IA e Dados\\[3pt]

    Belo Horizonte, MG
    \quad | \quad
    (37) 98815-1628
    \quad | \quad
    \href{mailto:amaralthalesalves@gmail.com}{amaralthalesalves@gmail.com}\\[2pt]

    \href{https://thalesamaral.tech}{thalesamaral.tech}
    \quad | \quad
    \href{https://www.linkedin.com/in/amaralthalesalves/}{linkedin.com/in/amaralthalesalves}
    \quad | \quad
    \href{https://github.com/amaralthalesalves}{github.com/amaralthalesalves}
\end{center}

% ================================================================
\section*{Resumo}

Engenheiro de software em P\&D no SPEED Lab (DCC/UFMG) e estudante de Ciência da Computação na UFMG.
Experiência com agentes de IA, LLMs (RAG, MCP, tool calling), pipelines de dados e observabilidade em Azure.

% ================================================================
\section*{Experiência Profissional}

\entry
{Engenheiro de Software (P\&D)}
{Nov. 2026 -- Presente}
{SPEED Lab -- Departamento de Ciência da Computação (DCC), UFMG}
{Belo Horizonte, MG}

\begin{itemize}
    \item Pesquisa e desenvolvimento de software com bolsa Embrapii.
\end{itemize}

\entry
{Estagiário em Engenharia de Dados e IA}
{Out. 2025 -- Out. 2026}
{Rehagro}
{Belo Horizonte, MG}

\begin{itemize}

    \item Automatizei a \textbf{transcrição e classificação de mais de 1.000 vídeos
    educacionais} com um pipeline em \textbf{Python, OpenAI Whisper e Apache Airflow},
    integrado aos bancos \textbf{PostgreSQL} internos.

    \item Implementei o \textbf{monitoramento de cerca de 80 máquinas virtuais Azure} com
    \textbf{Grafana e Prometheus}, apoiando análises de consumo, desempenho e custos de infraestrutura.

    \item Tornei dados de \textbf{VMs, custos e rede consultáveis em linguagem natural} com um \textbf{agente LLM com tool calling}.

    \item Estruturei o \textbf{catálogo de metadados} dos principais bancos corporativos,
    documentando regras de negócio, granularidade e cardinalidade, em uma aplicação
    \textbf{React/FastAPI integrada a um agente de IA}.

    \item Desenvolvi soluções com \textbf{RAG, MCP e avaliação de modelos}, dos requisitos com as áreas de negócio à validação.

    \item Integrei \textbf{agentes de IA, automações e APIs REST} à plataforma interna de Dados com \textbf{Laravel e React}.

    \item Antes do estágio: Treinamento em Engenharia de Software na Rehagro (jun. -- set. 2025).

\end{itemize}

\entry
{Técnico de Suporte em TI}
{Out. 2023 -- Set. 2024}
{Departamento de Toxicologia -- Universidade Federal de São João del-Rei (UFSJ)}
{Divinópolis, MG}

\begin{itemize}

    \item Mantive a \textbf{infraestrutura de TI de pesquisadores e laboratórios}, solucionando problemas de software, hardware e rede.

\end{itemize}

% ================================================================
\section*{Formação Acadêmica}

\entry
{Bacharelado em Ciência da Computação}
{Mar. 2025 -- 2029 (previsão)}
{Universidade Federal de Minas Gerais (UFMG)}
{Belo Horizonte, MG}

\begin{itemize}
    \item Disciplinas: Estruturas de Dados, Algoritmos, Bancos de Dados, Matemática Discreta e Probabilidade.
\end{itemize}

\entry
{Técnico em Informática Integrado ao Ensino Médio}
{Mar. 2022 -- Dez. 2024}
{Centro Federal de Educação Tecnológica de Minas Gerais (CEFET-MG)}
{Divinópolis, MG}

\begin{itemize}
    \item Desenvolvimento de software, bancos de dados, redes e sistemas computacionais. Representante de classe em 2022.
\end{itemize}

% ================================================================
\section*{Pesquisa e Projetos}

\entry
{Pesquisador Júnior -- Análise de Dados de Apreensões de Drogas Ilícitas}
{Jan. 2023 -- Dez. 2024}
{CEFET-MG, em parceria com UFSJ e Polícia Militar de Minas Gerais (PMMG)}
{Divinópolis, MG}

\begin{itemize}

    \item Estruturei e analisei dados de \textbf{apreensões de drogas, sistema prisional e informações
    geoespaciais} em \textbf{Python}, resultando em \textbf{artigo publicado}. Bolsa de Iniciação Científica
    \textbf{FAPEMIG} (out. 2023 -- set. 2024).

\end{itemize}

\entry
{FruitAdvisor -- Trabalho de Conclusão de Curso}
{2024}
{CEFET-MG}
{Divinópolis, MG}

\begin{itemize}

    \item Aplicativo mobile de \textbf{reconhecimento de frutas por imagem} com \textbf{Flutter} e uma \textbf{rede neural convolucional (CNN)}, exibindo informações nutricionais.

\end{itemize}

% ================================================================
\section*{Competências Técnicas}

\textbf{IA e Dados:}
LLMs, RAG, MCP, Tool Calling, Avaliação de Modelos, Machine Learning, OpenAI API, Whisper,
Airflow, PostgreSQL, MySQL\\
\textbf{Backend e Web:}
FastAPI, React, Node.js, Laravel, APIs REST\\
\textbf{Infraestrutura e Ferramentas:}
Azure, Docker, Grafana, Prometheus, Git, Linux, Shell\\
\textbf{Linguagens:}
Python, C, C++, Java, JavaScript, Dart, SQL, HTML, CSS

% ================================================================
\section*{Idiomas, Cursos e Prêmios}

\textbf{Machine Learning Specialization} -- DeepLearning.AI e Stanford Online (Coursera), Out. 2026 -- em andamento\\
\textbf{Inglês C1 (proficiência profissional plena)} -- Cambridge English Certification, pontuação \textbf{181/190}; Português nativo\\
\textbf{Medalha de Ouro} -- Olimpíada Nacional de Ciências (ONC), 2022

\end{document}
